%% =====================================================================
%% Section 5: A Formal Trade-off Framework for Representation
%%            Selection (RQ3)
%% =====================================================================
\section{A Formal Trade-off Framework for Representation Selection}
\label{sec:tradeoff}

\noindent\textit{``The choice of context representation is not an
implementation detail—it is the primary design decision that caps
a system's quality ceiling.''} --- derived from the SWE-bench analysis
of Yang et al.~\cite{yang2024sweagent}

\medskip

Sections~\ref{sec:taxonomy} and~\ref{sec:results} established \emph{what}
the seven PAR paradigms do and \emph{how well} they perform in controlled
benchmarks.  Section~\ref{sec:tradeoff} addresses the practitioner's
question: \emph{given a deployment context, which paradigm should I
choose?}  We answer this by constructing a \textbf{formal trade-off
framework} grounded in five operationally measurable dimensions, deriving
a scored comparison matrix, and stating explicitly when and why a
practitioner should prefer PAR-$X$ over PAR-$Y$.

% ----------------------------------------------------------------
\subsection{Motivation: Why One Dimension Is Never Enough}
\label{subsec:tradeoff:motivation}
% ----------------------------------------------------------------

Prior surveys present performance tables organised by benchmark
score~\cite{hou2024llmse_slr,husein2025llmcompletion_slr,tao2026ragsurvey}.
While informative, this single-axis ranking conflates fundamentally
different cost categories.  Consider two contrasting failure modes:

\begin{itemize}
  \item \textbf{Repoformer}~\cite{wu2024repoformer} achieves high
        token efficiency by selectively skipping retrieval for
        in-file-sufficient queries, but its static retrieval policy
        cannot adapt when a freshly pushed commit invalidates the
        indexed embedding space.

  \item \textbf{SWE-agent}~\cite{yang2024sweagent} dynamically
        discovers context by navigating the live repository, achieving
        near-perfect freshness, but issues 30--60 tool calls per task,
        incurring latency of 2--10 minutes on SWE-bench instances.
\end{itemize}

\noindent These failure modes are orthogonal: a system can be fast yet
stale, or fresh yet slow.  A practitioner deploying a CI/CD copilot
(\textit{e.g.}\  sub-second inline completion) faces entirely different
constraints than one deploying a nightly issue-resolution agent
(\textit{e.g.}\  SWE-bench-style).  The framework we introduce makes
these constraints explicit.

% ----------------------------------------------------------------
\subsection{The Five Operational Dimensions}
\label{subsec:tradeoff:dimensions}
% ----------------------------------------------------------------

We define five operational dimensions $\mathcal{D}_1$--$\mathcal{D}_5$
that jointly determine the cost--quality surface of a repository
representation paradigm.  Each dimension is assigned a score from
1 (worst) to 5 (best) based on evidence from primary studies.

\subsubsection*{$\mathcal{D}_1$: Semantic Precision}

\emph{Definition}: The fraction of task-relevant context that the
paradigm includes in the assembled prompt, excluding irrelevant noise.
Formally, let $R$ be the set of code elements causally required to solve
a task, and let $A$ be the set assembled by the paradigm.  Semantic
precision is $|R \cap A| / |A|$ (precision) and $|R \cap A| / |R|$
(recall); the harmonic mean $F_1$ serves as our proxy.

\emph{Evidence}: Graph-based systems (PAR-3) achieve the highest
$F_1$ on cross-file completion because dependency edges encode precise
causal links~\cite{liu2024graphcoder,cheng2024dataflow}.
GraphCoder's coarse-to-fine retrieval over code-context graphs
raises cross-file $F_1$ by $+$7.3 pp over BM25-based
RAG~\cite{liu2024graphcoder}.  PAR-2 systems with BM25 or dense
retrieval suffer precision degradation at scale: Gu et
al.~\cite{gu2025whatretrieve} show that \textit{what} is retrieved
matters more than \textit{how much}, and that token-budget saturation
(> 80\% of context filled with irrelevant chunks) cuts downstream
performance by 18\%.  PAR-7 agentic systems achieve high recall
by design (they can issue unlimited tool calls) but at the cost of
precision: the agent may navigate several red-herring files before
finding the root cause~\cite{yang2024sweagent,ma2025lingmaagent}.
PAR-6 compression introduces precision loss when the summariser
discards semantically thin but syntactically critical tokens
(\textit{e.g.}\  type annotations, exception signatures), a deficiency
documented by LongCodeZip~\cite{shi2025longcodezip} and quantified
by SWE-Pruner~\cite{wang2026swepruner}: over-aggressive pruning
incurs a $-$11\% task-completion regression on SWE-bench.

\subsubsection*{$\mathcal{D}_2$: Token Efficiency}

\emph{Definition}: The ratio of semantically informative tokens to
total tokens consumed, averaged over a task batch.  A higher ratio
implies more information per dollar of LLM API cost.

\emph{Evidence}: PAR-6 compression maximises token efficiency almost by
definition: LongCodeZip~\cite{shi2025longcodezip} achieves a 4$\times$
compression ratio (i.e.,  1{,}000 effective tokens from a 4{,}000-token
context) with $<$3\% downstream quality loss on code completion
benchmarks.  Repoformer~\cite{wu2024repoformer} selectively skips
retrieval on 35--55\% of queries that do not benefit from cross-file
context, saving one retrieval call and associated tokens per skipped
query with no quality regression.  PAR-3 graph systems incur a token
overhead because graph-serialisation formats (adjacency lists, Cypher
queries~\cite{liu2025codexgraph}) are verbose; CodexGraph's Cypher
interface reduces raw token usage vs.\  full-graph serialisation but
still costs 20--40\% more tokens than a pure BM25 RAG pipeline
for the same recall level.  PAR-7 agentic systems are worst: each tool
call consumes a full model forward pass plus the tool's output, and
multi-step navigation inflates total tokens by $3$--$8\times$ relative
to a single-shot PAR-2 retrieval~\cite{wang2025openhands}.  HumanEvo
\cite{zheng2025humanevo} further shows that token-budget waste in PAR-7
systems grows super-linearly with repository size, because agents explore
an increasing number of irrelevant file paths.

\subsubsection*{$\mathcal{D}_3$: Construction Latency}

\emph{Definition}: Wall-clock time to produce the context that is
submitted to the LLM, measured in a \emph{cold-start} scenario
(first query on a previously unseen repository).

\emph{Evidence}: PAR-1/PAR-2 systems with offline-indexed embeddings
achieve $<$100 ms retrieval latency after indexing, making them
suitable for interactive (sub-200 ms) code completion.  Repoformer
\cite{wu2024repoformer} adds a negligible classification overhead
($<$5 ms) to decide whether retrieval is needed at all.  PAR-3
systems incur indexing latency proportional to repository size:
building a project-level call graph for a 50K-LOC Java repository
takes 8--45 seconds depending on the tool
chain~\cite{liu2024graphcoder,ouyang2025repograph}.  At query time,
graph traversal is fast ($<$50 ms) for bounded-hop queries; however,
whole-graph operations for CodexGraph-style Cypher
queries~\cite{liu2025codexgraph} can reach 300--800 ms.  PAR-4
execution-trace systems depend on runtime: triggering a test suite
to obtain a stack trace takes 10s of seconds to minutes.
PAR-6 compression adds an amortisable one-time cost to
summarise repository files ($\sim$30--120 s for a 10K-LOC
codebase~\cite{shi2025longcodezip}), but the compressed representation
then serves any subsequent query instantly.  PAR-7 agentic navigation
is the highest-latency paradigm: median first-response time on
SWE-bench is 180--450 seconds for competitive
agents~\cite{yang2024sweagent,ma2025lingmaagent}.

\subsubsection*{$\mathcal{D}_4$: Update Cost}

\emph{Definition}: Incremental time and compute required to keep the
representation consistent after a single commit modifies a subset of
the repository.

\emph{Evidence}: PAR-2 RAG systems with chunk-level indices support
cheap \emph{incremental} updates: only modified file chunks need
re-embedding.  Repoformer~\cite{wu2024repoformer} tracks file
modification timestamps and re-indexes only changed chunks, keeping
update latency $O(\Delta)$ (proportional to diff size).  PAR-3
graph systems face a more expensive update: a commit that renames a
public API function invalidates all call-graph edges to that function,
potentially requiring whole-graph re-computation for languages without
incremental static-analysis
tooling~\cite{ouyang2025repograph,liu2024graphcoder}.  HumanEvo
\cite{zheng2025humanevo} stress-tests this by simulating repository
evolution across 12 months of real commit history; PAR-3 systems
degrade 14--22\% more than PAR-2 systems in the absence of
incremental graph updates.  CodeMEM~\cite{wang2026codemem} mitigates
update cost by representing repository state as a mutable in-context
memory tree that patches locally (AST node replacement) rather than
globally; update cost is $O(\text{diff lines})$.  PAR-7 agentic
systems have the lowest update cost: because they hold no persistent
index, they always read the live filesystem and are inherently
up-to-date, albeit at the cost of re-navigating from scratch on every
invocation~\cite{yang2024sweagent}.  PAR-6 compression systems must
re-summarise changed files; intelligent change-propagation strategies
are largely absent from the current literature (a gap noted in
\S\ref{subsec:tradeoff:gaps}).

\subsubsection*{$\mathcal{D}_5$: Scalability}

\emph{Definition}: Degradation in $\mathcal{D}_1$--$\mathcal{D}_4$
as repository size grows from 1K to 10K+ files.

\emph{Evidence}: PAR-1 dense-embedding indices scale
\emph{sub-linearly} in retrieval time thanks to ANNS (approximate
nearest-neighbour search) data structures, but their semantic coverage
degrades at scale because a fixed embedding dimension cannot
distinguish functionally similar modules~\cite{gu2025secret}.
SECRET~\cite{gu2025secret} addresses this with segmented deep hashing
and demonstrates $+$8\% retrieval quality at the 10K-file scale vs.\
standard FAISS indexing.  PAR-3 graph systems face $O(N^2)$ edge
growth for densely interconnected repositories; RepoHyper
\cite{phan2025repohyper} prunes the hypergraph to $k$-nearest
structural neighbours, maintaining $O(N \log N)$ index size.
PAR-7 agentic systems show the starkest scalability challenge:
Prometheus~\cite{pan2026prometheus} documents that agent trajectory
length grows linearly with codebase size, hitting context-window
limits ($\sim$200K tokens for GPT-4o) on repositories exceeding
$\sim$15K files.  Luo et al.'s RPG-Encoder~\cite{luo2026rpgencoder}
proposes universal repository embeddings trained to be robust across
scales, reporting stable performance from 500 to 50K files.
PAR-6 compression systems exhibit the best scalability profile
because the compressed representation is bounded by a fixed output
token count regardless of repository size~\cite{shi2025longcodezip}.

% ----------------------------------------------------------------
\subsection{Trade-off Matrix}
\label{subsec:tradeoff:matrix}
% ----------------------------------------------------------------

Table~\ref{tab:tradeoff} summarises the five-dimensional scoring for
each PAR paradigm.  Scores range from 1 (poor) to 5 (excellent)
and are assigned based on the evidence reviewed in
\S\ref{subsec:tradeoff:dimensions}.  An $\oplus$ marker indicates a
hybrid system that transcends its primary category.

\begin{table*}[t]
\caption{Trade-off Matrix: PAR Paradigms on Five Operational
Dimensions. Scores 1--5 (5~=~best). \textbf{Bold} indicates the
highest score per column. $\downarrow$ indicates lower is better for
$\mathcal{D}_3$ and $\mathcal{D}_4$ (we invert for display so 5 =
lowest latency / lowest update cost).}
\label{tab:tradeoff}
\small
\begin{tabular*}{\textwidth}{@{\extracolsep{\fill}}
    p{2.8cm}   % Paradigm
    >{\centering}p{1.7cm}  % D1
    >{\centering}p{1.7cm}  % D2
    >{\centering}p{1.7cm}  % D3
    >{\centering}p{1.7cm}  % D4
    >{\centering\arraybackslash}p{1.7cm}  % D5
  }
\toprule
\textbf{Paradigm} &
\textbf{$\mathcal{D}_1$}  &
\textbf{$\mathcal{D}_2$}  &
\textbf{$\mathcal{D}_3$}  &
\textbf{$\mathcal{D}_4$}  &
\textbf{$\mathcal{D}_5$}  \\
& \textit{Semantic} & \textit{Token} & \textit{Construct.} & \textit{Update} & \textit{Scala-} \\
& \textit{Precision} & \textit{Efficiency} & \textit{Latency}$^{\downarrow}$ & \textit{Cost}$^{\downarrow}$ & \textit{bility} \\
\midrule
PAR-1: Dense Embedding
  & 3 & 3 & \textbf{5} & 4 & 3 \\
\addlinespace[2pt]
PAR-2: RAG Cross-File
  & 3 & 4 & \textbf{5} & \textbf{5} & 4 \\
\addlinespace[2pt]
PAR-3: Graph-Based
  & \textbf{5} & 2 & 3 & 2 & 3 \\
\addlinespace[2pt]
PAR-4: Execution-Trace
  & 4 & 3 & 2 & \textbf{5} & 3 \\
\addlinespace[2pt]
PAR-5: Memory-Augmented
  & 4 & 3 & 4 & 4 & 3 \\
\addlinespace[2pt]
PAR-6: Compression
  & 2 & \textbf{5} & 4 & 3 & \textbf{5} \\
\addlinespace[2pt]
PAR-7: Agentic Navigation
  & 4 & 1 & 1 & \textbf{5} & 2 \\
\bottomrule
\end{tabular*}
\vspace{2pt}
\begin{minipage}{\textwidth}
\footnotesize
\textit{Evidence base}: $\mathcal{D}_1$: \cite{liu2024graphcoder,gu2025whatretrieve,wang2026swepruner};
$\mathcal{D}_2$: \cite{shi2025longcodezip,wu2024repoformer,liu2025codexgraph,wang2025openhands};
$\mathcal{D}_3$: \cite{wu2024repoformer,liu2024graphcoder,yang2024sweagent,shi2025longcodezip,ma2025lingmaagent};
$\mathcal{D}_4$: \cite{wu2024repoformer,zheng2025humanevo,wang2026codemem,yang2024sweagent};
$\mathcal{D}_5$: \cite{gu2025secret,phan2025repohyper,pan2026prometheus,shi2025longcodezip,luo2026rpgencoder}.
\end{minipage}
\end{table*}

Several patterns emerge from Table~\ref{tab:tradeoff}:

\begin{itemize}

  \item \textbf{No paradigm dominates.}  Each paradigm leads on at
        most two dimensions and trails on at least one.  The closest
        to an all-rounder is PAR-2, which scores $\geq 3$ on all
        dimensions—but it never leads on precision or scalability.

  \item \textbf{Precision and token efficiency are anti-correlated.}
        PAR-3 achieves the highest semantic precision (5) at the
        cost of the second-lowest token efficiency (2).  PAR-6
        achieves the highest token efficiency (5) at the cost of the
        lowest precision (2).  The fundamental tension is that rich
        structural information is verbose, and compression discards
        structural nuance.

  \item \textbf{Latency and update cost are positively correlated
        with precision for structural paradigms.}  PAR-3's high
        precision comes from expensive offline indexing (latency 3)
        and fragile incremental update (cost 2).  PAR-7's
        on-demand navigation achieves high freshness (update cost 5)
        at the price of catastrophic construction latency (1).

  \item \textbf{Scalability is an Achilles' heel for agentic
        systems.}  PAR-7 scores only 2 on scalability, the joint
        worst alongside PAR-3 (when graph size is unbounded).
        This is a critical constraint for industrial deployment.

\end{itemize}

% ----------------------------------------------------------------
\subsection{Practitioner Selection Guidelines}
\label{subsec:tradeoff:guidelines}
% ----------------------------------------------------------------

We now operationalise the matrix into concrete \emph{decision rules}
for system designers.  We frame each guideline as a conditional
statement on the deployment context.

\subsubsection{Guideline~G1: Choose PAR-2 when latency and update
freshness are hard constraints.}

Interactive tools such as IDE autocompletion have a perceptual latency
budget of $\leq$200~ms~\cite{wu2024repoformer}.  PAR-2 with offline
chunk-level indexing is the only paradigm that consistently satisfies
this budget after the cold-start indexing phase.  Its incremental
update model (re-embed only changed chunks) sustains freshness across
rapid commit cycles.  Repoformer's selective-retrieval
optimisation~\cite{wu2024repoformer} further reduces token waste for
in-file-sufficient queries.  \textbf{Practical implication}: a
production code-completion copilot should default to PAR-2; the
selective-retrieval gate of Repoformer should be treated as a
mandatory engineering best practice rather than an optional add-on.

\paragraph{Why not PAR-1?}  PAR-1 shares PAR-2's latency advantage
but lacks query-conditioned retrieval; it assembles context based on
the query embedding, missing cross-file dependencies that lie in
semantically distant but structurally adjacent files.
GraphCodeBERT~\cite{guo2021graphcodebert} partially bridges this gap
via DFG pre-training, but it remains a retrieval encoder, not a
context assembler.  \textbf{Verdict}: prefer PAR-2 over PAR-1 for
any completion task requiring cross-file context.

\subsubsection{Guideline~G2: Choose PAR-3 when downstream task quality
is the primary objective and repository churn is low.}

When token budget is generous, commit rate is low (\textit{e.g.}\  
a released library rather than an active development branch), and
task quality is paramount (\textit{e.g.}\  a production code-review
assistant), PAR-3 graph-based representation offers the highest
semantic precision.  GraphCoder~\cite{liu2024graphcoder}'s
coarse-to-fine retrieval over code-context graphs consistently
outperforms PAR-2 baselines by 5--9 pp on RepoBench's cross-file
completion categories.  RepoHyper~\cite{phan2025repohyper}'s
graph-pruning strategy limits memory growth on large repositories,
making PAR-3 feasible at moderate scale.

\paragraph{Why not PAR-3 for high-churn repositories?}
HumanEvo~\cite{zheng2025humanevo} quantifies the staleness penalty:
a PAR-3 system evaluated on 6-month-evolved repositories (without
graph refresh) degrades by 14--22\% relative to a contemporaneously
evaluated PAR-2 system.  Since call-graph refresh after a significant
API refactor is a non-trivial engineering operation, PAR-3 becomes
risky in rapid-release contexts.  \textbf{Verdict}: prefer PAR-3 for
static or slow-moving codebases; fall back to PAR-2 otherwise.

\subsubsection{Guideline~G3: Choose PAR-6 when operating under extreme
token budget constraints or serving very large repositories as a
single prompt.}

Large-context LLMs (128K--1M token windows) create a temptation to
simply concatenate all repository files.  In practice, input-length
degradation (``lost in the middle'' effects) and inference cost make
raw concatenation infeasible above $\sim$10K files.  PAR-6
compression provides a principled solution: LongCodeZip
\cite{shi2025longcodezip} achieves a 4$\times$ compression factor
with $<$3\% quality loss, enabling a 10K-LOC repository to fit in a
32K-token window.  Hierarchical summarisation
\cite{dhulshette2025hiersum} further adapts compression depth by
directory tier.

\paragraph{Why not PAR-6 for tasks requiring structural precision?}
LongCodeZip's training objective minimises cross-entropy on code
tokens, which biases compression toward \emph{frequent} patterns.
Rare but critical patterns—custom exception hierarchies, obscure
type constraints, deprecated-but-still-called APIs—are
disproportionately likely to be elided~\cite{wang2026swepruner}.
SWE-Pruner documents a $-$11\% regression on SWE-bench when
aggressive compression removes issue-relevant type
annotations~\cite{wang2026swepruner}.  \textbf{Verdict}: use PAR-6
when summarisation-robust tasks (documentation generation, high-level
planning) dominate; combine PAR-2 with PAR-6 for a hybrid that
retrieves precisely then compresses aggressively.

\subsubsection{Guideline~G4: Choose PAR-7 when the task requires
open-ended exploration and latency is not a constraint.}

SWE-bench~\cite{jimenez2024swebench} represents the class of tasks
where neither the relevant file set nor the root-cause code locus is
known in advance.  PAR-7 agentic systems are the only paradigm that
systematically outperforms all others on this task class, achieving
35--50\% resolution rates versus $<$18\% for PAR-2
\cite{ma2025lingmaagent,yang2024sweagent}.  The key insight is that
\emph{exploration is semantically necessary}: the agent must determine
which file to read before a meaningful context can be assembled.
No offline index can replicate this capability for arbitrary, unseen
issues.

\paragraph{Why not PAR-7 for large repositories?}
Prometheus~\cite{pan2026prometheus} identifies a \emph{horizon
collapse} failure mode: agentic trajectories for issues in
repositories with $>$15K files frequently exhaust the LLM's context
window on navigation scaffolding before reaching the root-cause file.
LingmaAgent~\cite{ma2025lingmaagent} mitigates this via a two-stage
coarse-to-fine localization; however, even this system degrades by
18\% on the largest SWE-bench repositories.  \textbf{Verdict}: use
PAR-7 unconditionally for complex issue-resolution tasks on
repositories up to $\sim$10K files; deploy PAR-7 with a
coarse-localisation pre-filter (PAR-3 or PAR-2) for larger
repositories.

\subsubsection{Guideline~G5: Prefer PAR-5 for iterative tasks where
context must accumulate across multiple LLM invocations.}

Tasks such as incremental feature development~\cite{wang2026codemem},
multi-turn refactoring, or long-horizon test generation
require the representation to \emph{grow} as the agent produces
and observes code changes.  PAR-5 memory systems maintain a
mutable, task-scoped representation that absorbs each edit
delta without discarding prior context.  CodeMEM
\cite{wang2026codemem} demonstrates a $+$9.1\% improvement over
PAR-2 on multi-turn code generation by preserving AST-node-level
change history in memory.  \textbf{Verdict}: for single-turn
tasks PAR-5 offers no advantage over PAR-2; activate PAR-5
only when inter-turn dependency is structurally necessary.

% ----------------------------------------------------------------
\subsection{Hybrid and Compositional Strategies}
\label{subsec:tradeoff:hybrid}
% ----------------------------------------------------------------

Real-world deployments increasingly combine paradigms rather than
applying a single one.  We identify three architecturally motivated
hybrid strategies:

\begin{description}
  \item[H1: PAR-2 $+$ PAR-3 (precision-guided retrieval).]
        Use a lightweight BM25/dense retriever (PAR-2) as a first
        stage to identify candidate files, then rerank and augment
        via structural graph traversal (PAR-3).  GraphCoder
        \cite{liu2024graphcoder} and RepoHyper
        \cite{phan2025repohyper} instantiate this strategy; both
        outperform pure PAR-2 and pure PAR-3 at the same token
        budget.

  \item[H2: PAR-7 $+$ PAR-3 (agent-guided structural navigation).]
        An agentic system (PAR-7) uses graph-structured indices
        (PAR-3) as a high-bandwidth navigation oracle rather than
        raw file reads.  LocAgent~\cite{chen2025sweexp} implements a
        PAR-7 outer loop that issues Cypher queries against a
        CodexGraph~\cite{liu2025codexgraph} database, reducing
        median tool calls from 48 to 21 on SWE-bench (a $-$56\%
        latency reduction) while retaining resolution quality.

  \item[H3: PAR-2 $+$ PAR-6 (efficient large-context retrieval).]
        Retrieve the top-$k$ chunks via PAR-2, then compress the
        assembled context via PAR-6 before submitting to the LLM.
        This strategy is applicable when token budget is tight but
        quality must exceed PAR-6 alone.  LongCodeZip
        \cite{shi2025longcodezip} can be applied in this post-hoc
        compression role without modification; the 4$\times$
        compression rate is largely preserved even on
        RAG-assembled contexts.
\end{description}

% ----------------------------------------------------------------
\subsection{Open Problems and Research Gaps}
\label{subsec:tradeoff:gaps}
% ----------------------------------------------------------------

The trade-off framework exposes several open problems that the
current literature does not adequately address:

\begin{enumerate}

  \item \textbf{Incremental PAR-3 graph maintenance.}
        No existing PAR-3 system provides a production-grade
        incremental graph-update protocol that is robust to API
        renames, file moves, and dependency removals.  HumanEvo
        \cite{zheng2025humanevo} quantifies the staleness penalty
        but does not provide a remedy.  A formal change-propagation
        algorithm for code-context graphs—analogous to incremental
        compilers in language tooling—is an open engineering problem.

  \item \textbf{Precision-preserving compression for PAR-6.}
        Current compression models~\cite{shi2025longcodezip} are
        trained to minimise generation loss, which is a poor proxy
        for task-relevant information retention.  A compression
        objective that explicitly penalises the elision of
        structurally critical tokens (type annotations, exception
        signatures, API contracts) is absent from the literature.

  \item \textbf{Scalable PAR-7 trajectory control.}
        Prometheus~\cite{pan2026prometheus} documents horizon
        collapse but does not solve it.  A principled algorithm for
        \emph{adaptive trajectory truncation}—terminating navigation
        when information gain per tool call falls below a threshold—
        would substantially improve PAR-7 scalability.

  \item \textbf{Multi-project and cross-repository representation.}
        All seven PAR paradigms assume a single-repository context.
        Microservice architectures~\cite{jiang2026surveycodegen} and
        monorepos with loosely coupled sub-packages violate this
        assumption.  The interaction between inter-repository
        dependency graphs and intra-repository context representation
        is unexplored.

  \item \textbf{Unified evaluation protocol across dimensions.}
        Benchmark suites including SWE-bench~\cite{jimenez2024swebench},
        RepoBench~\cite{liu2024repobench}, and HumanEvo
        \cite{zheng2025humanevo} each capture one or two of the five
        dimensions (primarily $\mathcal{D}_1$) but none jointly
        measures all five.  A community-standard benchmark that
        reports latency, update cost, and scalability alongside
        task quality would substantially accelerate comparative
        research.

\end{enumerate}

% ----------------------------------------------------------------
\subsection{Summary of RQ3}
\label{subsec:tradeoff:summary}
% ----------------------------------------------------------------

\begin{formal}
\noindent\textbf{Answer to RQ3.}
No single repository representation paradigm is optimal across all
five operational dimensions.  PAR-2 (RAG cross-file) provides the
best all-round baseline, particularly for latency-constrained
interactive tools.  PAR-3 (graph-based) maximises semantic precision
for stable repositories but is fragile under high commit rates.  PAR-6
(compression) is the most token-efficient and scalable paradigm but
sacrifices structural precision.  PAR-7 (agentic navigation) is the
only paradigm that reliably resolves open-ended repository issues, but
its latency and scalability limitations restrict its applicability to
batch or asynchronous deployment contexts.  Practitioners operating
near the Pareto frontier of the cost--quality surface should adopt
hybrid strategy H1 (PAR-2 $+$ PAR-3) for structured completion tasks
and H2 (PAR-7 $+$ PAR-3) for issue-resolution tasks, reserving
pure PAR-6 for large-repository summarisation pipelines.
\end{formal}
